\ifintegriereaufgaben

\section*{Laplace Transformation Anwendungen Aufgaben}\label{MA2-laplacetransformation-und-differentialgleichungen-aufgaben}
Erkl�rungen und Beispiele, die zur Bearbeitung der Aufgaben hilfreich sind, finden Sie im Skript\wennreferenzvorhanden{MA2-laplacetransformation-und-differentialgleichungen}{ (Abschnitt \abschnittsverweis{MA2-laplacetransformation-und-differentialgleichungen})} und, falls Sie nicht weiterkommen, dann schauen Sie \hyperref[MA2-laplacetransformation-und-differentialgleichungen-loesungen]{hier}.
\else
\section{Laplace Transformation Anwendungen Aufgaben}\label{MA2-laplacetransformation-und-differentialgleichungen-aufgaben}

Regeln und Beispiele zu diesen Aufgaben finden Sie im Skript. Die L�sungen finden Sie \hyperref[MA2-laplacetransformation-und-differentialgleichungen-loesungen]{hier}.
\fi

\begin{exercisebox}[Laplace-Transformation Rechenregeln] % Papula 2, VI, Abschnitt 2, Aufgabe 1, 2, 3, 5, 6, 7, 9, 10
	Bestimmen Sie durch Anwendung der Regeln die Laplace-Transformierte der angegebenen Funktionen:
	\begin{itemize}
		\item[(a)] $f(t) = 4t^3-t^2+2t $
		\item[(b)] $f(t) = C(1-e^{-\lambda t}) $
		\item[(c)] $f(t) = \left( 3t\right)^5 $ Hinweis: $t^5 \,\laplace\, \displaystyle{\frac{120}{s^6}} $
		\item[(d)] $f(t) = \cos( 4t) $ Hinweis: $\cos(t) \,\laplace\, \displaystyle{\frac{s}{s^2 +1}} $
		\item[(e)] $f(t) = \cos^2( \omega t) $ Hinweis: $\cos^2(t) \,\laplace\, \displaystyle{\frac{s^2+2}{s\cdot (s^2 +4)}} $
			\item[(f)] $f(t) = (t-4)^2 $ Hinweis: $t^2 \,\laplace\, \displaystyle{\frac{2}{s^3}} $
			\item[(g)] $f(t) = e^{t-b}, \; b>0 $ Hinweis: $e^t \,\laplace\, \displaystyle{\frac{1}{s-1}} $
			\item[(h)] $f(t) = \sin(t+\pi/2) $ Hinweis: $\sin(t) \,\laplace\, \displaystyle{\frac{1}{s^2+1}} $
			\item[(i)] $f(t) = e^{3t}\cdot \cos(2t) $ Hinweis: $\cos(t) \,\laplace\, \displaystyle{\frac{s}{s^2+1}} $
			\item[(j)] $f(t) = A\cdot e^{-\delta t}\cdot \sin(\omega t), \; \delta, \omega >0 $ Hinweis: $\sin(t) \,\laplace\, \displaystyle{\frac{1}{s^2+1}} $
			\item[(k)] $f(t) = 2^{3t}$ Hinweis: $e^{\lambda t} \,\laplace\, \displaystyle{\frac{1}{s-\lambda}} $
		\end{itemize}
\end{exercisebox}


\begin{exercisebox}[Laplace-Transformation und Differentialgleichungen]
	Betrachten Sie die Differentialgleichung $ y'' + \omega^2 y = 0,\; \omega>0$ mit den beiden Anfangswerten
	\begin{enumerate}
		\item $y(0) = 1 \land y'(0) = 0$
		\item $y(0) = 0 \land y'(0) = 1$
	\end{enumerate}
	\begin{itemize}
		\item[(a)] Berechnen Sie die L�sung der beiden Anfangswertprobleme.
		\item[(b)] Berechnen Sie die Laplace-Transformierte $Y(s)$ der Differentialgleichung \( y'' + \omega^2\, y = 0\). Nutzen Sie dabei die Ableitungsregeln.
		\item[(c)] Bestimmen Sie, ausgehend von den speziellen L�sungen der Anfangswertprobleme im Bildbereich, die L�sungen im Zeitbereich. Verwenden Sie f�r die R�cktransformationen eine Tabelle.
	\end{itemize}
\end{exercisebox}

\begin{exercisebox}[Laplace-Transformation und Differentialgleichungen]
	L�sen Sie die folgenden Systeme mit Hilfe der Laplace-Transformation
	\begin{eqnarray*}
		(a)\; \begin{cases}
			x(0) = y(0) =  0, \quad & t= 0\,,\\
			\dot x + 2 y = e^t, \quad & t> 0\,,\\
			\dot y + 2 x = e^{-t}, \quad & t> 0\,.
		\end{cases} \quad\quad
		(b)\; \begin{cases}
			x(0) = y(0) =  0, \quad & t= 0\,,\\
			\dot x(0) = \dot y(0) =  0, \quad & t= 0\,,\\
			\ddot x + \dot y + 3 y = 1, \quad & t> 0\,,\\
			\ddot y -4 \dot x + 3y = 0, \quad & t> 0\,.
		\end{cases} \quad\quad
	\end{eqnarray*}
\end{exercisebox}


%\begin{exercisebox}[Laplace-Transformation und Differentialgleichungen]  %HMI III; Klausur, Aufgabe 2
%	Die Analyse eines elektromagnetischen Reihenschwingkreises habe Sie zu folgender Integro-Differentialgleichung f�r die Stromst�rke $I(t)$ gef�hrt:
%	\[
%	\dot I(t) + 2\delta I(t) + \omega_0^2 \int\limits_0^t I(\tau)\,\mathrm{d}\tau = \frac{U_0}{L} \land I(0)=0\,.\]
%	Hierbei sind $\delta, \omega_0, L$ und $U_0$ positive Konstanten und es gelte $\omega^2_0 - \delta^2 >0$. L�sen Sie dieses Anfangswertproblem mit Hilfe der Laplace-Transformation. Hinweis: F�r die R�cktransformation k�nnen Sie die folgende Formel ohne Beweis verwenden:
%	\[ \mathcal L^{-1} \left( \frac{1}{(s+b)^2 + a^2}\right) = \frac{d^{-bt}\sin(at)}{a}, \; a>0\,.\]
%\vspace{0.3cm}
%{\bf L�sung:} \newline
%\end{exercisebox}
